% 由 generate-clustering-data.py 生成的真实 OPTICS 排序教学例。
% m=3（含自身），eps_max=1.1；无穷值用顶部三角形显示。
\begin{tikzpicture}
 \begin{axis}[
   width=\linewidth,height=27mm,at={(0,48mm)},anchor=south west,
   axis x line=bottom,axis y line=none,
   xmin=-0.4,xmax=12.5,ymin=-0.5,ymax=0.5,
   xtick={0,4,8,12},ytick=\empty,
   xlabel={原始特征值 $x$},font=\small,
   axis line style={FigureMuted},tick label style={FigureMuted}]
  \addplot[only marks,mark=*,mark size=1.8pt,FigureInk]
   table[x=x,y expr=0] {figures/03-models/01-classic-models/tikz/clustering-optics.dat};
 \end{axis}
 \begin{axis}[
   width=\linewidth,height=47mm,at={(0,0)},anchor=south west,
   xmin=0.4,xmax=23.6,ymin=0,ymax=1.15,
   axis lines=left,xlabel={OPTICS 访问次序},ylabel={距离},
   xtick={1,8,14,23},ytick={0,0.4,0.8},
   font=\small,axis line style={FigureMuted},
   tick label style={FigureMuted},unbounded coords=jump,
   legend style={draw=none,font=\footnotesize,at={(0.5,-0.42)},
     anchor=north,legend columns=3}]
  \path[fill=FigureBlueFill] (axis cs:.5,0) rectangle (axis cs:7.5,1.15);
  \path[fill=FigureYellowFill] (axis cs:13.5,0) rectangle (axis cs:22.5,1.15);
  \addplot[ycomb,FigureRed,mark=*,mark size=1.5pt,line width=1.1pt]
   table[x=order,y=reach] {figures/03-models/01-classic-models/tikz/clustering-optics.dat};
  \addlegendentry{可达距离}
  \addplot[FigureBlue,dashed,line width=1pt]
   table[x=order,y=core] {figures/03-models/01-classic-models/tikz/clustering-optics.dat};
  \addlegendentry{核心距离}
  \addplot[only marks,mark=triangle*,mark size=2.6pt,FigureStroke]
   table[x=order,y=restart] {figures/03-models/01-classic-models/tikz/clustering-optics.dat};
  \addlegendentry{可达距离$+\infty$}
  \draw[FigureStroke,dash dot,line width=1pt]
    (axis cs:0.5,0.35) -- (axis cs:23.5,0.35);
  \node[text=FigureInk,anchor=south,inner sep=1pt]
    at (axis cs:18,0.40) {$\varepsilon'=0.35$};
 \end{axis}
\end{tikzpicture}
